\documentclass[11pt,a4paper,fontset=fandol]{ctexart}

\usepackage[a4paper,top=2.4cm,bottom=2.6cm,left=2.35cm,right=2.35cm,headheight=18pt,headsep=0.55cm,footskip=26pt]{geometry}
\usepackage{amsmath,amssymb}
\usepackage{fancyhdr}
\usepackage{lastpage}
\usepackage{enumitem}
\usepackage{titlesec}
\usepackage{xcolor}
\usepackage{hyperref}
\usepackage{bookmark}
\usepackage[most]{tcolorbox}
\usepackage{setspace}
\usepackage{array}

\hypersetup{
  colorlinks=true,
  linkcolor=blue!50!black,
  urlcolor=blue!50!black
}

\setstretch{1.16}
\setlength{\parindent}{2em}
\setlength{\parskip}{0.32em}
\setlist[enumerate]{itemsep=0.35em, topsep=0.35em, leftmargin=2.3em}
\setlist[itemize]{itemsep=0.25em, topsep=0.25em, leftmargin=2em}

\definecolor{mainblue}{RGB}{36,83,140}
\definecolor{lightblue}{RGB}{241,246,252}
\definecolor{lightgray}{RGB}{246,246,246}
\definecolor{rulegray}{RGB}{110,118,128}
\definecolor{softgreen}{RGB}{236,247,240}
\definecolor{softyellow}{RGB}{254,249,229}

\titleformat{\section}
  {\Large\bfseries\color{mainblue}}
  {\thesection.}
  {0.45em}
  {}
  [\vspace{0.25em}\color{rulegray}\titlerule]
\titlespacing*{\section}{0pt}{1.2em}{0.8em}
\titleformat{\subsection}{\large\bfseries\color{mainblue}}{\thesubsection}{0.5em}{}

\pagestyle{fancy}
\fancyhf{}
\fancyhead[L]{\small 图论计数专题目录页}
\fancyhead[C]{\small 组合专题资料索引}
\fancyhead[R]{\small 刘文博}
\fancyfoot[C]{\small 第 \thepage 页 / 共 \pageref{LastPage} 页}
\renewcommand{\headrulewidth}{0.55pt}
\renewcommand{\footrulewidth}{0.35pt}

\newtcolorbox{goalbox}[1]{
  breakable,
  colback=lightblue,
  colframe=mainblue,
  boxrule=0.7pt,
  arc=1mm,
  title=\textbf{#1},
  fonttitle=\bfseries
}

\newtcolorbox{infobox}[1]{
  breakable,
  colback=softgreen,
  colframe=green!40!black,
  boxrule=0.65pt,
  arc=1mm,
  title=\textbf{#1},
  fonttitle=\bfseries
}

\newtcolorbox{warnbox}[1]{
  breakable,
  colback=softyellow,
  colframe=orange!60!black,
  boxrule=0.65pt,
  arc=1mm,
  title=\textbf{#1},
  fonttitle=\bfseries
}

\begin{document}

\thispagestyle{empty}
\begin{titlepage}
\centering
\vspace*{0.9cm}

\begin{tcolorbox}[
  enhanced,
  width=0.92\textwidth,
  colback=white,
  colframe=mainblue,
  boxrule=1pt,
  arc=0mm,
  left=6mm,
  right=6mm,
  top=8mm,
  bottom=8mm
]
\centering

{\fontsize{24}{30}\selectfont\bfseries 图论计数专题目录页\par}
\vspace{0.65cm}
{\fontsize{17}{22}\selectfont 讲义、题单与周测的使用说明\par}

\vspace{1cm}
\begin{tcolorbox}[colback=lightblue,colframe=mainblue,width=0.84\textwidth,boxrule=1pt]
\centering
{\large 适合对象：已经掌握基础计数，准备系统提升“把关系翻译成图”的学生}\\[0.35em]
{\normalsize 本目录页帮助把图论计数专题资料串成一个完整训练包}
\end{tcolorbox}

\vspace{1.0cm}
\renewcommand{\arraystretch}{1.42}
\begin{tabular}{>{\bfseries}p{3.5cm}p{8cm}}
专题名称： & 图论入门中的组合计数 \\
资料组成： & 课堂讲义、提高训练题单、周测 A 卷、周测 B 卷 \\
使用方式： & 先学讲义，再做题单，最后用 A/B 周测检查掌握情况 \\
编译方式： & XeLaTeX \\
\end{tabular}

\vspace{0.9cm}
{\large \today\par}
\end{tcolorbox}
\end{titlepage}

\tableofcontents
\newpage

\section{专题包包含哪些资料}

\begin{goalbox}{四份核心资料}
\begin{enumerate}
  \item 课堂讲义：帮助理解点、边、度数、握手定理、完全图、二分图与网格路径；
  \item 提高训练题单：用于课后巩固和变式提升，重点训练“会建图、会数边、会数路径、会做结构证明”；
  \item 周测 A 卷：检查基础图论建模与基本计数是否稳定；
  \item 周测 B 卷：检查结构证明、网格路径变式和经典图论结论的理解是否真正到位。
\end{enumerate}
\end{goalbox}

\begin{infobox}{对应文件名}
\begin{itemize}
  \item 讲义：\texttt{graph\_counting\_handout.tex} / \texttt{graph\_counting\_handout.pdf}
  \item 提高训练题单：\texttt{graph\_counting\_advanced\_problem\_set.tex} / \texttt{graph\_counting\_advanced\_problem\_set.pdf}
  \item 周测 A 卷：\texttt{graph\_counting\_weekly\_test\_A.tex} / \texttt{graph\_counting\_weekly\_test\_A.pdf}
  \item 周测 B 卷：\texttt{graph\_counting\_weekly\_test\_B.tex} / \texttt{graph\_counting\_weekly\_test\_B.pdf}
\end{itemize}
\end{infobox}

\section{建议学习顺序}

\begin{goalbox}{推荐顺序}
\begin{enumerate}
  \item 先完整学习讲义，重点吃透“什么是点、边、度数”“为什么握手定理成立”“网格路径为什么能转化成计数”；
  \item 再完成提高训练题单，训练把题目翻译成完全图、二分图或网格图；
  \item 接着做周测 A 卷，检查基础模型是否稳定；
  \item 最后做周测 B 卷，检查结构证明和综合变式是否掌握。
\end{enumerate}
\end{goalbox}

\begin{warnbox}{不建议的做法}
\begin{itemize}
  \item 只会套公式，不先判断题目对应哪一种图；
  \item 用握手定理时忘了“每条边被数两次”；
  \item 网格路径题只会画图，不会转成组合数；
  \item 做证明题时不画图、不分类，直接凭感觉写结论。
\end{itemize}
\end{warnbox}

\section{每份资料主要解决什么问题}

\subsection{课堂讲义}
讲义适合第一次系统学习这个专题，主要解决以下问题：
\begin{itemize}
  \item 如何把“人、城市、格点、比赛关系”翻译成图；
  \item 握手定理为什么成立；
  \item 完全图和二分图的边数怎么数；
  \item 网格路径为什么本质上是图上的路径计数；
  \item 一些经典结构题如何用图论语言重新表达和证明。
\end{itemize}

\subsection{提高训练题单}
题单适合第二阶段训练，主要用来把方法做熟。它更强调：
\begin{itemize}
  \item 能否正确建图；
  \item 能否在边数、度数和路径数之间切换；
  \item 能否处理“经过某点”“避开某点”的路径变式；
  \item 能否对命题先判断真假，再给出证明或反例。
\end{itemize}

\subsection{周测 A 卷}
A 卷偏基础，适合检查：
\begin{itemize}
  \item 握手定理是否熟练；
  \item 完全图、二分图边数是否会数；
  \item 基础网格路径是否会转化为组合数；
  \item 简单聚会、比赛问题是否会建图。
\end{itemize}

\subsection{周测 B 卷}
B 卷偏综合，适合检查：
\begin{itemize}
  \item 是否会做图论中的结构证明；
  \item 是否会处理网格路径中的“经过/避开”点；
  \item 是否会识别并反驳错误命题；
  \item 是否能把抽屉原理、染色与图论语言结合起来。
\end{itemize}

\section{一个推荐的 4 次训练安排}

\begin{goalbox}{4 次完成一个小专题}
\begin{enumerate}
  \item 第 1 次：学习讲义前半部分，重点理解点、边、度数与握手定理；
  \item 第 2 次：学习讲义后半部分，完成题单基础题和中档题；
  \item 第 3 次：完成题单综合题，并把错题按“建图错了/公式错了/路径分类错了/命题判断错了”分类订正；
  \item 第 4 次：完成周测 A 卷或 B 卷，作为阶段性检查。
\end{enumerate}
\end{goalbox}

\begin{infobox}{如果孩子基础较好，可以这样压缩}
\begin{itemize}
  \item 第 1 天：讲义 + 题单基础题；
  \item 第 2 天：题单变式题 + 错题订正；
  \item 第 3 天：周测 A 卷；
  \item 第 4 天：周测 B 卷。
\end{itemize}
\end{infobox}

\section{专题学习后的自检清单}

\begin{goalbox}{如果下面 8 条都能做到，说明这个专题基本过关}
\begin{enumerate}
  \item 我能准确说出点、边、度数分别表示什么；
  \item 我会使用握手定理；
  \item 我会计算完全图的边数；
  \item 我会计算二分图的最大边数；
  \item 我会把网格路径题转化成组合数；
  \item 我会处理“经过某点”或“避开某点”的路径变式；
  \item 我能用图论语言重新表述一些经典结论；
  \item 我会在证明题里先判断命题真假，再决定是证明还是举反例。
\end{enumerate}
\end{goalbox}

\section{给家长的使用建议}

\begin{warnbox}{更有效的带练方式}
\begin{itemize}
  \item 不要一开始就讲公式，先让孩子说“点代表谁、边代表什么”；
  \item 如果卡住，可以只提示“能不能把它画成图”或“能不能数度数和”；
  \item 网格路径题订正时，重点看孩子是否真正理解“走法对应排列”；
  \item 结构证明题建议先让孩子尝试画图，再比较不同方法。
\end{itemize}
\end{warnbox}

\begin{infobox}{和后续专题的衔接}
图论计数学完以后，可以自然过渡到图染色、欧拉路径、哈密顿路径、图论中的抽屉原理与极值问题等更进一步的竞赛专题。
\end{infobox}

\end{document}
